%
%Academic CV LaTeX Template
% Author: Dubasi Pavan Kumar
% LinkedIn: https://www.linkedin.com/in/im-pavankumar/
% License: MIT
%
% For errors, suggestions, or improvements, please contact:
% Email: pavankumard.pg19.ma@nitp.ac.in
%



\documentclass[a4paper,11pt]{article}

% Package imports
\usepackage{graphicx}

\usepackage{latexsym}
\usepackage{xcolor}
\usepackage{float}
\usepackage{ragged2e}
\usepackage[empty]{fullpage}
\usepackage{wrapfig}
\usepackage{lipsum}
\usepackage{tabularx}
\usepackage{titlesec}
\usepackage{geometry}
\usepackage{marvosym}
\usepackage{verbatim}
\usepackage{enumitem}
\usepackage{fancyhdr}
\usepackage{multicol}
\usepackage{graphicx}
\usepackage{cfr-lm}
\usepackage[T1]{fontenc}
\usepackage{fontawesome5}
\usepackage{enumitem}

% Color definitions
\definecolor{darkblue}{RGB}{0,0,139}

% Page layout
\setlength{\multicolsep}{0pt} 
\pagestyle{fancy}
\fancyhf{} % clear all header and footer fields
\fancyfoot{}
\renewcommand{\headrulewidth}{0pt}
\renewcommand{\footrulewidth}{0pt}
\geometry{left=1.4cm, top=0.8cm, right=1.2cm, bottom=1cm}
\setlength{\footskip}{5pt} % Addressing fancyhdr warning

% Hyperlink setup (moved after fancyhdr to address warning)
\usepackage[hidelinks]{hyperref}
\hypersetup{
    colorlinks=true,
    linkcolor=darkblue,
    filecolor=darkblue,
    urlcolor=darkblue,
}

% Custom box settings
\usepackage[most]{tcolorbox}
\tcbset{
    frame code={},
    center title,
    left=0pt,
    right=0pt,
    top=0pt,
    bottom=0pt,
    colback=gray!20,
    colframe=white,
    width=\dimexpr\textwidth\relax,
    enlarge left by=-2mm,
    boxsep=4pt,
    arc=0pt,outer arc=0pt,
}

% URL style
\urlstyle{same}

% Text alignment
\raggedright
\setlength{\tabcolsep}{0in}

% Section formatting
\titleformat{\section}{
  \vspace{-4pt}\scshape\raggedright\large
}{}{0em}{}[\color{black}\titlerule \vspace{-7pt}]

% Custom commands
\newcommand{\resumeItem}[2]{
  \item{
    \textbf{#1}{\hspace{0.5mm}#2 \vspace{-0.5mm}}
  }
}

\newcommand{\resumePOR}[3]{
\vspace{0.5mm}\item
    \begin{tabular*}{0.97\textwidth}[t]{l@{\extracolsep{\fill}}r}
        \textbf{#1}\hspace{0.3mm}#2 & \textit{\small{#3}} 
    \end{tabular*}
    \vspace{-2mm}
}

\newcommand{\resumeSubheading}[4]{
\vspace{0.5mm}\item
    \begin{tabular*}{0.98\textwidth}[t]{l@{\extracolsep{\fill}}r}
        \textbf{#1} & \textit{\footnotesize{#4}} \\
        \textit{\footnotesize{#3}} &  \footnotesize{#2}\\
    \end{tabular*}
    \vspace{-2.4mm}
}

\newcommand{\resumeProject}[4]{
\vspace{0.5mm}\item
    \begin{tabular*}{0.98\textwidth}[t]{l@{\extracolsep{\fill}}r}
        \textbf{#1} & \textit{\footnotesize{#3}} \\
        \footnotesize{\textit{#2}} & \footnotesize{#4}
    \end{tabular*}
    \vspace{-2.4mm}
}

\newcommand{\resumeSubItem}[2]{\resumeItem{#1}{#2}\vspace{-4pt}}

\renewcommand{\labelitemi}{$\vcenter{\hbox{\tiny$\bullet$}}$}
\renewcommand{\labelitemii}{$\vcenter{\hbox{\tiny$\circ$}}$}

\newcommand{\resumeSubHeadingListStart}{\begin{itemize}[leftmargin=*,labelsep=1mm]}
\newcommand{\resumeHeadingSkillStart}{\begin{itemize}[leftmargin=*,itemsep=1.7mm, rightmargin=2ex]}
\newcommand{\resumeItemListStart}{\begin{itemize}[leftmargin=*,labelsep=1mm,itemsep=0.5mm]}

\newcommand{\resumeSubHeadingListEnd}{\end{itemize}\vspace{2mm}}
\newcommand{\resumeHeadingSkillEnd}{\end{itemize}\vspace{-2mm}}
\newcommand{\resumeItemListEnd}{\end{itemize}\vspace{-2mm}}
\newcommand{\cvsection}[1]{%
\vspace{2mm}
\begin{tcolorbox}
    \textbf{\large #1}
\end{tcolorbox}
    \vspace{-4mm}
}

\newcolumntype{L}{>{\raggedright\arraybackslash}X}%
\newcolumntype{R}{>{\raggedleft\arraybackslash}X}%
\newcolumntype{C}{>{\centering\arraybackslash}X}%

% Commands for icon sizing and positioning
\newcommand{\socialicon}[1]{\raisebox{-0.05em}{\resizebox{!}{1em}{#1}}}
\newcommand{\ieeeicon}[1]{\raisebox{-0.3em}{\resizebox{!}{1.3em}{#1}}}

% Font options
\newcommand{\headerfonti}{\fontfamily{phv}\selectfont} % Helvetica-like (similar to Arial/Calibri)
\newcommand{\headerfontii}{\fontfamily{ptm}\selectfont} % Times-like (similar to Times New Roman)
\newcommand{\headerfontiii}{\fontfamily{ppl}\selectfont} % Palatino (elegant serif)
\newcommand{\headerfontiv}{\fontfamily{pbk}\selectfont} % Bookman (readable serif)
\newcommand{\headerfontv}{\fontfamily{pag}\selectfont} % Avant Garde-like (similar to Trebuchet MS)
\newcommand{\headerfontvi}{\fontfamily{cmss}\selectfont} % Computer Modern Sans Serif
\newcommand{\headerfontvii}{\fontfamily{qhv}\selectfont} % Quasi-Helvetica (another Arial/Calibri alternative)
\newcommand{\headerfontviii}{\fontfamily{qpl}\selectfont} % Quasi-Palatino (another elegant serif option)
\newcommand{\headerfontix}{\fontfamily{qtm}\selectfont} % Quasi-Times (another Times New Roman alternative)
\newcommand{\headerfontx}{\fontfamily{bch}\selectfont} % Charter (clean serif font)

\begin{document}
\headerfontiii

% ===== Header with photo on right =====
% ===== Header with photo (clean 2-line contacts) =====
% ===== Header with photo: tighter + aligned for resume =====
% \noindent
% \begin{minipage}[t]{0.78\textwidth}
%     \vspace{2mm} % pushes name+links DOWN (increase/decrease)
%     {\Huge\textbf{Avni Mittal}}\\[-1.0mm]

%     \small
%     +91-9667730877 \;|\;
%     \href{mailto:avni.mittal2002@gmail.com}{avni.mittal2002@gmail.com} \;|\;
%     \href{https://scholar.google.com/citations?user=itURCtMAAAAJ&hl=en}{Google Scholar}
%     \\[-0.5mm]
%     \socialicon{\faLinkedin}\;
%     \href{https://linkedin.com/in/avni-mittal-5a0b3b198/}{avni-mittal-5a0b3b198}
%     \;|\;
%     \socialicon{\faGithub}\;
%     \href{https://avnimittal13.github.io/portfolio/}{AvniMittal13}
% \end{minipage}%
% \hfill
% \begin{minipage}[t]{0.20\textwidth}
%     \raggedleft
%     \vspace{-5mm} % pulls photo UP (more negative = higher)
%     \includegraphics[width=2.35cm]{passport_photo.jpeg} % smaller photo
% \end{minipage}


% \vspace{-3mm}
% \hrule
% \vspace{-5mm}


% Header
\begin{center}
    {\Huge\textbf{Avni Mittal}}
\end{center}
\vspace{-4mm}

\begin{center}
    \small{
    +91-9667730877 | \href{mailto:avni.mittal2002@gmail.com}{avni.mittal2002@gmail.com} 
    |
    \href{https://scholar.google.com/citations?user=itURCtMAAAAJ&hl=en}{Google Scholar} 
    |
    \socialicon{\faLinkedin} \href{linkedin.com/in/avni-mittal-5a0b3b198/}{avni-mittal-5a0b3b198} | 
    \socialicon{\faGithub} \href{https://avnimittal13.github.io/portfolio/}{AvniMittal13} 
    % | 
    % \href{https://www.yourwebsite.com/}{yourwebsite.com}
    }
\end{center}
\vspace{-5mm}


\section{\textbf{Education}}
\vspace{0.5mm}
\resumeSubHeadingListStart

\resumeSubheading
  {Indian Institute of Technology Mandi}{Himachal Pradesh, India}
  {Bachelor of Technology in Computer Science and Engineering (Honours) - with Minor in Management}{Nov. 2020 -- May 2024}
  \resumeItemListStart
    \item \small{CGPA: 8.93/10, Department Rank 2}

  \resumeItemListEnd

\resumeSubheading
  {Technische Universität Darmstadt}{Darmstadt, Germany}
  {Semester Exchange (5th Semester), B.Tech CSE}{Oct. 2022 -- Mar. 2023}
  \resumeItemListStart
    \item \small{CGPA: 9.36/10}

  \resumeItemListEnd

\resumeSubHeadingListEnd
\vspace{-6.5mm}

\section{\textbf{Research Interests}}
\vspace{0.5mm}
\small{Agentic misalignment, AI alignment and safety, interpretability, and reasoning in LLMs. I am interested in understanding model internals and failure cases, then using those insights to build lightweight, interpretable, and accessible AI systems inspired by the brain's energy efficiency.}

\vspace{-2mm}

\section{\textbf{Ongoing Projects}}
\vspace{0.5mm}
\resumeItemListStart
  \item \textbf{TextNCA} - Neural cellular automata for language modeling.
  \item \textbf{sonar-moe} - Speculative routing of tokens in Mixture-of-Experts (MoE) models.
  \item \textbf{VIAS} - Organizational value alignment via activation steering.
  \item \textbf{Understanding Cross-Image Information Flow in Multi-Image Vision-Language Models}.
\resumeItemListEnd

\vspace{-2mm}

\section{\textbf{Patents and Publications} \hfill \textcolor{darkblue}{\scriptsize C=Conference, P=Patent}}
\vspace{0.5mm}
\small{
\begin{enumerate}[leftmargin=*, labelsep=0.5em, align=left, widest={[\textbf{S.1}]}, itemindent=0em, label={\textbf{[\arabic*]}}]

% ----------------- CONFERENCES -----------------

% -------------------- PUBLICATIONS --------------------

% \item[\textbf{[C.1]}] 
% \textit{\underline{\textbf{Avni Mittal}}}, Arka Mukhopadhyay, Arnav Kulkarni, Anhad Thukral, Ayuj Aryan, Arnav Bhavsar.  
% “Grapher-NCA: Integrating Global Context via Graph-Based Updates in Neural Cellular Automata.”  
% [In Preparation]  
% \href{}{Link (to be added)}

\item[\textbf{[C.1]}] 
\textit{\underline{\textbf{Avni Mittal}}}, Shanu Kumar, Sandipan Dandapat, Monojit Choudhary.  
“LITMUS++: An Agentic System for Predictive Analysis of Low-Resource Languages Across Tasks and Models.”  
\textit{IJCNLP–AACL 2025, System Demonstration Track}.  
\href{https://aclanthology.org/2025.ijcnlp-demo.6/}{[Link]}

\item[\textbf{[C.2]}] 
\textit{\underline{\textbf{Avni Mittal}}}, Sree Hari Nagaralu, Sandipan Dandapat.  
“PROTECT: Policy-Related Organizational Value Taxonomy for Ethical Compliance and Trust.”  
\textit{SICon 2025 (ACL Workshop)}.  
\href{https://aclanthology.org/2025.sicon-1.5/}{[Link]}

\item[\textbf{[C.3]}] 
\textit{\underline{\textbf{Avni Mittal}}}, John Kalkhof, Anirban Mukhopadhyay, Arnav Bhavsar.  
“MedSegDiffNCA: Diffusion Models with Neural Cellular Automata for Skin Lesion Segmentation.”  
\textit{IEEE CBMS 2025}.  
\href{https://www.computer.org/csdl/proceedings-article/cbms/2025/261000a035/284Khy3rT3i}{[Link]}

\item[\textbf{[C.4]}] 
Avinash Anand, \textit{\underline{\textbf{Avni Mittal}}}, Laavanaya Dhawan, Mahisha Ramesh, Juhi Krishnamurthy,  
Naman Lal, Raj Jaiswal, Pijush Bhuyan, Himani, Astha Verma, Rajiv Ratn Shah, Roger Zimmermann, Shin’ichi Satoh.  
“Unveiling Learner Dynamics: The ECLIPSE Dataset and NeuralGaze Framework for Prolonged Engagement Assessment in Online Learning.”  
\textit{ECAI 2024}.  
\href{https://ieeexplore.ieee.org/abstract/document/10707841}{[Link]}

\item[\textbf{[C.5]}] 
Avinash Anand, \textit{\underline{\textbf{Avni Mittal}}}, Laavanaya Dhawan, Juhi Krishnamurthy, Mahisha Ramesh,  
Naman Lal, Astha Verma, Pijush Bhuyan, Himani, Rajiv Ratn Shah, Roger Zimmermann, Shin’ichi Satoh.  
“Exceda: Unlocking Attention Paradigms in Extended Duration E-Classrooms with Attention Mechanism.”  
\textit{IEEE MIPR 2024}.  
\href{https://ebooks.iospress.nl/doi/10.3233/FAIA240551}{[Link]}


% -------------------- PATENTS --------------------

\item[\textbf{[P.1]}] 
\textit{\underline{\textbf{Avni Mittal}}}, Shanu Kumar, Sandipan Dandapat, Monojit Choudhary.  
“Multi-Agent Predictive Systems for Cross-Domain Model Performance Assessment”  
\textit{Patent filed}.

\item[\textbf{[P.2]}] 
Gautam Prasad, \textit{\underline{\textbf{Avni Mittal}}}, Yugal Sachdev.  
“A System of Simultaneous Optimization of Multiple Agents.”  
\textit{Patent filed}.

\item[\textbf{[P.3]}] 
Gautam Prasad, \textit{\underline{\textbf{Avni Mittal}}}, Yugal Sachdev.  
“A System to Navigate Product Complexities through Automatic Visual Guide Generation.”  
\textit{Patent being filed}.

\end{enumerate}
}

\vspace{-2mm}

\section{\textbf{Research Experience}}
\vspace{0.5mm}

\resumeSubHeadingListStart

\resumeSubheading
  {Multi-Agent System for Predictive Analysis of Low-Resource Languages}{July 2025 – Present}
  {Supervisors: Dr. Sandipan Dandapat and Dr. Monojit Choudhury}{}
  \resumeItemListStart
 \item Designed a multi-agent system based on directed acyclic graph (DAG) guided by an expert knowledge graph to predict model performance for the task language without fine-tuning or full benchmarking; agents retrieve multilingual evidence, extract relevant features and train lightweight regressors for prediction.
    \item Created benchmark scenarios across six tasks and five multilingual settings to evaluate predictive accuracy and reasoning quality under diverse low-resource conditions.
    \item Work accepted at \textbf{IJCNLP-AACL 2025 (System Demonstration Track)}; patent filed [\textbf{C.1}, \textbf{P.1}].
  \resumeItemListEnd

\resumeSubheading
  {Faithfulness and Deceptive Reasoning in Mixture-of-Experts Models}{Sept. 2025 – Present}
  {Supervised Research at SPAR AI (Safety Alignment Research)}{}
  \resumeItemListStart
    % \item Studying faithfulness and deceptive reasoning in Mixture-of-Experts (MoE) models by analyzing expert-level token routing and evaluating how disabling high-deception experts affects reasoning faithfulness and task accuracy.
    % \item Developing a unified faithfulness metric combining accuracy and reasoning alignment with token level importance based decomposition, and extending it through \textbf{Thought Anchors} for cross-domain reasoning evaluation.

% \item Designing cue-based benchmarks to test whether each step in a model’s chain-of-thought causally follows from prior reasoning and injected cues, and developing a unified faithfulness metric combining causal faithfulness and completeness.
\item Developed a perturbation-based benchmark from PRM800K to evaluate Chain-of-Thought (CoT) faithfulness using causality and completeness as core criteria to benchmark LLMs-as-judges for faithfulness detection.

% \item Designing another cue-based benchmark and unified faithfulness metrics to evaluate whether each step in a model’s chain-of-thought causally follows from prior reasoning and injected cues.

    \item Studying faithfulness and deceptive reasoning in Mixture-of-Experts (MoE) models by analyzing expert-level token routing and evaluating how disabling high-deception experts affects reasoning faithfulness and task accuracy.


  \resumeItemListEnd

\resumeSubheading
  {Medical Image Segmentation using Neural Cellular Automata and Diffusion}{Aug. 2023 – May. 2024}
  {Supervisors: Dr. Arnav Bhavsar and Dr. Anirban Mukhopadhyay}{}
  \resumeItemListStart
 \item Conducted bachelor's thesis research on enhancing the \textbf{M3D-NCA} framework for medical image segmentation through region-of-interest detection and Vision Transformer–based Cellular Automata (ViTCA), achieving \textasciitilde40\% faster inference with comparable accuracy.
    \item Proposed the integration of diffusion models with Neural Cellular Automata for improved segmentation; published as \textbf{MedSegDiffNCA} at \textbf{IEEE CBMS 2025} [\textbf{C.3}].
\item Currently extending this work through \textbf{Grapher-NCA}, integrating graph-based global reasoning into NCA architectures to improve segmentation accuracy and convergence, while mentoring an undergraduate student team on model design and experimentation.


  \resumeItemListEnd

\resumeSubheading
  {Multimodal Engagement Prediction in Classroom Environments}{Sept. 2023 – Aug. 2024}
  {Supervisor: Dr. Rajiv Ratn Shah, MIDAS Lab IIIT Delhi, Research Intern}{}
  \resumeItemListStart
  \item Developed novel pipeline for multimodal engagement prediction in classroom environments using Neural Cellular Automata (NCA) based framework integrating local-global attention and behavioural facial features.
  \item Curated large-scale multimodal datasets (\textbf{ExCEDA}, \textbf{ECLIPSE}) for prolonged engagement and behavioral analysis.
  % \item Proposed a \textbf{Neural Cellular Automata (NCA)} based framework integrating local-global attention and behavioral facial features, achieving improved efficiency and performance over existing baselines.
    \item Introduced a Context-Guided Vision Transformer (ViT) approach for modeling temporal shifts in engagement, capturing transitions across four affective states—boredom, engagement, frustration, and confusion.
  \item Research outcomes accepted at \textbf{ECAI 2024} and \textbf{IEEE MIPR 2024} \textbf{[C.4, C.5]}.
\resumeItemListEnd

% \resumeSubheading
%   {ToolBot: Tool Learning for Large Language Models}{Oct. 2023 -- Dec. 2023}
%   {Independent Research Project}{}
%   \resumeItemListStart
%   \item Designed a conversational agent (\textbf{ToolBot}) for developer assistance using enterprise tool-specific knowledge and context-aware LLM retrieval; achieved \textbf{5th position at Inter-IIT Tech Meet 12.0}.

  % \resumeItemListEnd

\resumeSubheading
  {Explainable AI Methods and Robustness Analysis}{Dec. 2022 -- Mar. 2023}
  {Student Research Assistant, TU Darmstadt}{}
\resumeItemListStart
  \item Conducted a comprehensive literature review of explainable AI (XAI) methods and implemented approaches including Integrated Gradients, Grad-CAM, and SmoothGrad.
  \item Evaluated robustness of XAI methods under adversarial attacks (DeepFool, Carlini–Wagner L2) by analyzing changes in focus maps to measure sensitivity of XAI visualizations under perturbations.
\resumeItemListEnd

\resumeProject
  {C-Arm Orientation Tracking for Surgical Workflow Optimization}{Dec. 2022 -- Mar. 2023}
  {Dr. Anirban Mukhopadhyay, TU Darmstadt)}{}
\resumeItemListStart
\item Developed real-time methods for 3D C-arm orientation estimation from 2D images using ML and analytical geometry approaches; complete pipeline implemented and preliminary work submitted to \textbf{MICCAI-AE CAI 2023}.

\resumeItemListEnd


\resumeSubheading
  {Human Behavior Modeling and Yoga Pose Classification}{Dec. 2021 -- Mar. 2022}
  {Research Intern, with Dr. Varun Dutt and Dr. Arnav Bhavsar, IIT Mandi iHub \& HCL}{}
\resumeItemListStart
  \item Developed LSTM and LRCN-based pose classification method for real-time gamified environments.
  \item Modeled human gameplay behavior using LSTM-based next-action prediction to analyze strategic decision-making and behavioral patterns in reinforcement learning settings for Unity-RL agents.
  \item Work recognized among top 10 nationwide teams at the \textit{ACM-W Hackathon 2021}.
\resumeItemListEnd

\resumeSubHeadingListEnd

\vspace{-3mm}
\section{\textbf{Industry Experience}}
\vspace{0.5mm}

\resumeSubHeadingListStart

\resumeSubheading
{Microsoft}{Hyderabad, India}
{Data Scientist (AI Research)}{Jan. 2024 -- Present}
\textit{Microsoft Security Platform - Research in multi-agent systems, AI alignment, and compliance automation.}
\resumeItemListStart
  % \item Conducting applied research on multi-agent architectures, AI safety alignment, and organizational value compliance frameworks.
  \item Developed the \textbf{Security Analysis Agent}, evolving from chained workflows to autonomous agentic designs with context engineering and self-correcting reasoning, achieving \textbf{89\%} precision and earning the \textit{AI Innovation Award}. Currently working to improve security reasoning depth and consistency (planning-coding) through RL.
  \item Proposed and implemented the \textbf{Multi-Agent Prompt Tuner}, a scalable framework for coordinated prompt optimization across interdependent agents using hierarchical prompt design, textual gradient-based feedback, and dynamic reward generation to enable simultaneous tuning while maintaining task coherence (\textit{patent filed} \textbf{[P.2]}).  \item Designed the \textbf{Swarm Pinboard}, a multi-agent orchestration framework \textit{inspired by swarm intelligence}, enabling autonomous collaboration across a graph of intelligent nodes with reasoning capabilities for complex workflows.
  % \item Built a \textbf{Security Copilot Skill} for natural language–driven data-table search and API generation. Optimized payload generation and API mapper modules; conducted extensive accuracy testing. Improved RAG performance from \textbf{60\% to 91\%}.
  \item Built a \textbf{Security Copilot Skill} for natural language–driven data-table search and API generation, deployed to over 2,000 enterprise tenants. Optimized payload generation and API mapper modules; conducted extensive accuracy testing. Improved RAG performance from \textbf{60\% to 91\%}.

\item Built an efficient multi-label document classification and reasoning for Sensitive Information Type (SIT) detection using BitNet Transformer, combining chunk-embedding encoder with language-based decoder.
  
  \item Developed a multi-agent system for visual guide generation for enterprise workflows (\textit{patent being filed} \textbf{[P.3]}).
  \item Proposed a policy-driven value taxonomy integrating human ethics and compliance objectives for assessing LLM organizational alignment, resulting in \textbf{ACL SICon 2025} publication \textbf{[C.2]}.
\resumeItemListEnd

\resumeSubheading
{Intuit India}{Bangalore, India}
{Software Engineer Intern}{May 2023 -- Jul. 2023}
\resumeItemListStart
  \item Redesigned the internal \textit{TestEasy} platform from a monolithic to a microservices-based architecture, integrating downstream document, consent and identity pipelines.
  \item Implemented and deployed a production-ready test account cloning feature with unit and integration tests.
\resumeItemListEnd

\resumeSubheading
{Fellowship.AI}{Remote}
{Data Science Fellow}{Oct. 2022 -- Dec. 2022}
\resumeItemListStart
  \item Reproduced and modularized the \textbf{Total Relighting} pipeline for portrait relighting and background replacement.
  \item Integrated open-source models (Semantic Guided Human Matting, GracoNet Object Placement, Harmonizer) into an interactive real-time demo with seamless pipeline orchestration for end-to-end processing.
\resumeItemListEnd
\resumeSubHeadingListEnd

\vspace{-5mm}

\section{\textbf{Standardized Tests}}
\vspace{0.5mm}
\textbf{TOEFL (iBT):} \textbf{119/120} - Listening: 30, Reading: 29, Speaking: 30, Writing: 30\\
\textbf{GRE General Test:} \textbf{326/340} - Quant: 168, Verbal: 158, Analytical Writing: 4.0

\vspace{-2mm}

\section{\textbf{Awards and Talks}}
\vspace{0.5mm}
\resumeItemListStart
  \item \textbf{Invited Talk}, "How Agentic Systems Can Be Developed for Organizational Information Processing Using RAG." at IIT Hyderabad for Digital Transformation series \hfill \textit{\footnotesize{Sept. 2025}}
  
  \item \textbf{AI Innovation Award}, Microsoft Security Platform - For Security Analysis Agent \hfill \textit{\footnotesize{2025}}

  \item \textbf{Institute Merit Scholarship (Female Students)} - Full tuition waiver awarded for academic excellence. \hfill \textit{\footnotesize{2020-2024}}
  
  \item \textbf{Academic Excellence Award (CSE)}, IIT Mandi - Highest SGPA in 5th and 6th semesters. \hfill \textit{\footnotesize{2023, 2024}}
  
  \item \textbf{Top 5 Rank}, Inter-IIT Tech Meet - Representing IIT Mandi. Also achieved Top 10 Rank in 2022. \hfill \textit{\footnotesize{2023}}
  
  \item \textbf{Top 10 Finalist}, ACM-W India Hackathons \hfill \textit{\footnotesize{2021, 2022}}
  
  \item \textbf{Winner}, LLM Hack 2024, Atlas Astrothon, Vision Arcadia, Build-a-Bot, and Hack36 organized at IIT Mandi \hfill \textit{\footnotesize{2024}}

   \item \textbf{JEE Main AIR 5067} among 1.2M candidates; \textbf{JEE Advanced AIR 5287} among 250k candidates \hfill \textit{\footnotesize{2020}}

\resumeItemListEnd

\vspace{-2mm}

\section{\textbf{Teaching and Leadership Experience}}
\vspace{0.5mm}

\resumeItemListStart


\item \textbf{CS672 Advanced Deep Learning} \hfill Feb.-May 2025 \\
\textit{Research group lead, guided two UG student groups on GNNs and NCA, Advised by Dr. Arnav Bhavsar} \hfill \textit{SCEE, IIT Mandi}

\vspace{1.5mm}

\item \textbf{CS202 Paradigms of Programming} \hfill Feb.-May 2024 \\
\textit{Teaching assistant, advised by Dr. Varunkumar Jayapaul} \hfill \textit{SCEE, IIT Mandi}

\vspace{1.5mm}

\item \textbf{IC272 Data Science III} \hfill Aug.-Dec. 2023 \\
\textit{Teaching assistant, advised by Dr. Dilip A. D.} \hfill \textit{SCEE, IIT Mandi}

\item \textbf{Core Team Member - Conducting sessions and Organizing talks} \hfill 2021-2023 \\
\textit{ACM-Women Chapter, Robotronics Club, Space Technology and Astronomy Club, Programming Club} \hfill \textit{IIT Mandi}


\resumeItemListEnd

\vspace{-2mm}

\section{\textbf{Skills and Relevant Courses}}
\vspace{0.5mm}
\resumeHeadingSkillStart

  \resumeSubItem{Relevant Courses: }
    {Advanced Deep Learning, Computer Vision, Learning Robots, Probabilistic Graphical Models, Text Analytics, Information and Database systems, Computer Orgnanization, Data Structures and Algorithms}

% \resumeSubItem{Machine Learning \& Deep Learning: }
%     {Diffusion Models, Neural Cellular Automata (NCA), Vision Transformers, Large Language Models (LLMs), Prompt Engineering}

  \resumeSubItem{Programming Languages: }
    {Python, C++, Java, JavaScript, C, Shell Scripting, LaTeX}

  \resumeSubItem{Frameworks \& Tools: }
    {LangChain, AutoGen, LangGraph, FastAPI, Docker, Git, Azure, PyTorch, TensorFlow}


\resumeHeadingSkillEnd

\vspace{-2mm}

\section{\textbf{References}}
\vspace{0.5mm}

\resumeItemListStart

\item \textbf{Dr. Monojit Choudhury}\\
Professor of Natural Language Processing\\
Mohamed bin Zayed University of Artificial Intelligence. \\
% \textit{Research Interests:} Medical Image Computing, Computer-Assisted Surgery, Domain Adaptation \\
\textbf{Email:} monojit.choudhury@mbzuai.ac.ae

\item \textbf{Dr. Sandipan Dandapat}\\
Associate Professor, Department of Computer Science and Engineering\\
Indian Institute of Technology Hyderabad \\
% \textit{Research Interests:} Natural Language Processing, Generative AI, Multilingual Language Technologies \\
\textbf{Email:} sdandapat@cse.iith.ac.in

\item \textbf{Dr. Rajiv Ratn Shah}\\
Professor, MIDAS Lab\\
Indraprastha Institute of Information Technology Delhi \\
% \textit{Research Interests:} Computer Vision, Medical Image Analysis, Deep Learning \\
\textbf{Email:} rajivratn@iiitd.ac.in

\item \textbf{Dr. Arnav Bhavsar}\\
Associate Professor, School of Computing and Electrical Engineering\\
Indian Institute of Technology Mandi \\
% \textit{Research Interests:} Computer Vision, Medical Image Analysis, Deep Learning \\
\textbf{Email:} arnav@iitmandi.ac.in

\item \textbf{Dr. Anirban Mukhopadhyay}\\
Independent Research Group (IRG) Leader, MEC Lab\\
Technische Universität Darmstadt, Germany \\
% \textit{Research Interests:} Medical Image Computing, Computer-Assisted Surgery, Domain Adaptation \\
\textbf{Email:} anirban.mukhopadhyay@gris.tu-darmstadt.de
\resumeItemListEnd

% \section{\textbf{References}}
% \vspace{-0.2mm}
% \small{
% \begin{enumerate}[leftmargin=*,labelsep=2mm]
% \item \textbf{Dr. Sandipan Dandapat}\\
%    Job Title, Department\\
%    Organization/Institution Name\\
%    Email: email1@example.com\\
%    Phone: +X-XXX-XXX-XXXX\\
%    \textit{Relationship: [e.g., Thesis Advisor, Manager, etc.]}

% \item \textbf{Dr. Arnav Vinayak Bhavsar}\\
%    Job Title, Department\\
%    Organization/Institution Name\\
%    Email: email2@example.com\\
%    Phone: +X-XXX-XXX-XXXX\\
%    \textit{Relationship: [e.g., Project Supervisor, Colleague, etc.]}

% \item \textbf{Dr. Anirban Mukhopadhyay}\\
%    Job Title, Department\\
%    Organization/Institution Name\\
%    Email: email3@example.com\\
%    Phone: +X-XXX-XXX-XXXX\\
%    \textit{Relationship: [e.g., Mentor, Collaborator, etc.]}
% \end{enumerate}
% }

% \vspace{-2mm}
% \section{\textbf{Links}}
% \vspace{0.5mm}
% Google Scholar: \url{https://scholar.google.com/citations?user=itURCtMAAAAJ&hl=en}\\
% LinkedIn: \url{https://www.linkedin.com/in/avni-mittal-5a0b3b198/}\\
% Portfolio: \url{https://avnimittal13.github.io/portfolio/}

% \vspace{1em}
% \noindent\textbf{Publication Links:}\\[0.25em]

% \textbf{C.2} PROTECT (SICon 2025): \url{https://aclanthology.org/2025.sicon-1.5/}\\

% \textbf{C.3} MedSegDiffNCA (IEEE CBMS 2025): \url{https://www.computer.org/csdl/proceedings-article/cbms/2025/261000a035/284Khy3rT3i}\\
% \textbf{C.4} ECLIPSE Dataset (ECAI 2024): \url{https://ieeexplore.ieee.org/abstract/document/10707841}\\
% \textbf{C.5} Exceda (IEEE MIPR 2024): \url{https://ebooks.iospress.nl/doi/10.3233/FAIA240551}





\end{document}